\section{Исходное сравнение пяти реализаций и промежуточная диагностика}
\label{app:historical-hardware}
\subsection{Время вычисления и согласованность повторного выполнения}
\begin{table}[t]
\caption{Исходный протокол с пятью реализациями: измеренное время в мкс, среднее из пяти средних по проходам из 500 потоков и выборочное стандартное отклонение этих пяти средних. Средние по проходам на Mega совпали.}
\label{tab:latency}
\centering\footnotesize
\begin{tabular}{lrr}
\toprule
Реализация & Mega 2560 & ESP32-C3\\
\midrule
KAN с коэффициентами & $1396.568\pm0$ & $5.474\pm0.113$\\
KAN с LUT отсчётов & $243.040\pm0$ & $5.889\pm0.201$\\
Многослойная KAN & $12306.176\pm0$ & $55.351\pm0.257$\\
MLP16 & $1328.672\pm0$ & $19.180\pm0.142$\\
DT5 & $17.400\pm0$ & $1.420\pm0.269$\\
\bottomrule
\end{tabular}
\end{table}

\begin{figure}[!tp]
\centering\includegraphics[width=\textwidth]{hardware_latency.pdf}
\caption{Время вычислительного ядра на физических устройствах для одних и тех же 500 исходных потоков в каждом проходе. По горизонтали используется логарифмическая шкала времени; отрезки погрешности показывают выборочное стандартное отклонение пяти средних по проходам, разделённым сбросами. Тактовые частоты плат, программные платформы и наборы команд различаются. Отсутствие наблюдаемого разброса между повторами на Mega не означает нулевую неопределённость измерения.}
\label{fig:latency}
\end{figure}

Таблица~\ref{tab:latency} и рис.~\ref{fig:latency} показывают, что при исходном протоколе с пятью реализациями соотношение времени и объёма хранения при выборе представления зависит от процессора. На Mega LUT уменьшает среднее время в 5.746 раза относительно KAN с коэффициентами. На ESP32-C3 среднее составляет 5.8888~мкс против 5.4740~мкс для коэффициентов: наблюдается увеличение на 0.4148~мкс, или 7.58\%, а не ускорение. Эксперимент не выделяет влияние кеша, USB или отдельных машинных команд, поэтому конкретное причинное объяснение этой разницы не устанавливается.

DT5 имеет наименьшее наблюдаемое среднее время на обоих устройствах. Многослойная KAN требует существенно больше времени, чем аддитивные представления, особенно на AVR. Эти наблюдения относятся к зафиксированным реализациям и не задают ранжирование архитектур независимо от их представления, качества предсказаний и предобработки.

Для каждой платы получены 25 записей и 12,500 измеренных решений. Все 25,000 решений совпадают с соответствующими скомпилированными эталонными реализациями на C; все 12,500 попарно соответствующих предсказаний совпадают между платами. Это согласованность повторного выполнения, а не безошибочная классификация. На 500 различных исходных потоках KAN с коэффициентами, KAN с LUT, многослойная KAN, MLP16 и DT5 допускают соответственно 18, 18, 6, 6 и 9 ошибок классификации. Два аддитивных представления дают одинаковые решения на всём общем наборе. Эти числа ошибок не заменяют метрики на полной отложенной тестовой выборке.

Отдельные значения времени лежат на сетке с шагом 4~мкс на Mega и 1~мкс на ESP32-C3. Приводятся исходные длительности без вычитания накладных расходов таймера. Пустые измерительные пробы на Mega обычно дают 0--4~мкс; в записи для многослойной модели диапазон составляет 0--8~мкс. На ESP диапазон составляет 0--2~мкс. Это особенно существенно для очень короткого ядра DT5. Одинаковые последовательности времени на Mega во всех повторах дают нулевое выборочное стандартное отклонение, но не означают нулевой неопределённости времени. Последовательности на ESP различаются. Использовано по одному экземпляру платы каждого типа и фиксированный порядок реализаций; стандартное отклонение характеризует вариацию повторов на этих платах, а не производственный разброс или независимые рабочие нагрузки.

\subsection{Отдельная диагностика размещения данных на ESP32-C3}
\label{sec:c3diagnostic}
Затем мы проверили, влияют ли размещение в памяти и повторное обращение к одному входу на соотношение затрат двух представлений на той же ESP32-C3. Этот отдельный эксперимент сохраняет зафиксированные веса, $L=1025$ и те же 500 исходных потоков. Одно приложение содержит два ядра, принимающих указатели на параметры: коэффициентное и табличное. Каждое ядро без изменений используется с параметрами во Flash или с их копией в DRAM, что даёт четыре условия. Ядра и измерительные функции размещены в IRAM; при компиляции сохранён флаг \texttt{-Os}, оптимизация на этапе компоновки отключена. Размещение проверено по секциям ELF и выведенным адресам. Две копии параметров занимают дополнительно 20,808 байт DRAM, помимо измерительных буферов и дескрипторов; это не измерение пикового расхода SRAM.

Пять упорядоченных проходов выполняются в пределах одной загрузки. Перед каждым из десяти блоков сбора выполняются 2,000 прогревочных вызовов, охватывающих все 500 потоков во всех четырёх условиях. Порядок по схеме Уильямса для четырёх условий уравновешивает позиции и попарное предшествование на всём наборе. С начала прогрева до конца каждого блока сбора операции последовательного интерфейса отсутствуют; прерывания остаются разрешёнными. Счётчиком тактов CPU измеряется либо один вызов с возвратом полного целочисленного логита (\texttt{stream}), либо 32 вызова на одном входе (\texttt{same32}). Последний интервал включает цикл и сохранение 32 результатов; деление на 32 даёт усреднённые затраты в составе блока, а не задержку отдельного вызова. Накладные расходы таймера не вычитаются. При подтверждённой частоте 160~МГц деление числа тактов на 160 даёт микросекунды. Измерение заканчивается возвратом целочисленного логита; сравнение с порогом, входившее в исходный протокол, исключено.

\begin{table}[t]
\caption{Отдельная диагностика ESP32-C3 в тактах на вызов: среднее и выборочное стандартное отклонение пяти средних по проходам в пределах одной загрузки. Столбец \texttt{same32} представляет усреднённые затраты блока из 32 вызовов, включая цикл и сохранение результатов.}
\label{tab:c3diagnostic}
\centering\small
\begin{tabular}{lrr}
\toprule
Параметры & \texttt{stream} & \texttt{same32}\\
\midrule
Коэффициенты, Flash & $881.41\pm5.96$ & $832.51\pm0.34$\\
LUT, Flash & $771.80\pm9.15$ & $501.41\pm0.29$\\
Коэффициенты, DRAM & $820.24\pm1.93$ & $830.39\pm0.26$\\
LUT, DRAM & $482.52\pm1.61$ & $492.96\pm0.43$\\
\bottomrule
\end{tabular}
\end{table}

В рамках этой диагностики LUT уменьшает средние затраты в режиме \texttt{stream} в 1.142 раза при размещении во Flash и в 1.700 раза при размещении в DRAM (табл.~\ref{tab:c3diagnostic}). Выигрыш среднего не распространяется на все характеристики задержки: 95-й процентиль по объединённым 2,500 вызовам \texttt{stream}, вычисленный линейной интерполяцией и округлённый до целых тактов, составляет 1,482 для LUT во Flash против 1,119 для коэффициентов во Flash. Перенос параметров из Flash в DRAM уменьшает это среднее на 37.48\% для LUT и на 6.94\% для коэффициентов. Для LUT разность Flash--DRAM составляет 289.28 такта на вызов в режиме \texttt{stream}, но 8.45 такта в среднем на вызов в составе блока \texttt{same32}. Следовательно, доступ к памяти и локальность входов существенно влияют на наблюдаемые затраты, а ускорение LUT на ESP32-C3 возможно. Счётчик промахов кэша не регистрировался. Эти наблюдения не устанавливают единственную причину исходного увеличения времени на 7.58\%: между экспериментами также изменились размещение кода, доступ к параметрам, таймер, граница измеряемого интервала и протокол последовательного интерфейса. Исходные измерения сохранены как отдельный результат.

Диагностический архив содержит 20,000 строк, представляющих 330,000 измеренных вызовов. Проверки в прошивке не выявили расхождений полных логитов с зафиксированными эталонами соответствующих представлений. Логиты одиночных вызовов передаются по отдельности; каждая строка блока из 32 вызовов содержит последний логит, контрольную сумму и число расхождений после проверки всех 32 результатов вне измеряемого интервала. Эти проверки устанавливают согласованность выполненных диагностических ядер с их эталонами; пять проходов в пределах одной загрузки не являются независимыми повторениями на разных платах или после отдельных загрузок.
